% !TeX root = ../main.tex
\section{Effect Recognition and Parameter Estimation}
\label{sec:recognition}

\subsection{Single effects on isolated notes}

Stein et al.\ detected single and cascaded effects on IDMT notes with
hand-crafted spectral, cepstral and harmonic features and support vector
machines, reaching 99.2\% for combinations of six effects on mono guitar
notes~\cite{stein2010aes,stein2010dafx}. Schmitt and Schuller showed that such
features transfer poorly across conditions: from monophonic to polyphonic
material, recall fell to chance with functionals and to 41--62\% with a
bag-of-audio-words model~\cite{schmitt2017recognising}. They also noted that
delay is trivial to detect on single notes and hard on phrases, where echoes
look like repeated notes. Jürgens et al.\ added parameter regression with
shallow networks, with mean absolute errors of 0.05--0.16 on a unit scale for
distortion, delay and tremolo controls~\cite{jurgens2020recognizing}.
Comunità et al.\ trained a CNN (FxNet) to identify 13 drive pedal plugins
(81--91\% accuracy) and a class-conditioned regressor for gain and tone (MAE
0.02--0.07)~\cite{comunita2021guitar}. Most confusions were between nearly
identical circuits such as the TS808 and TS9, a first sign that the inverse
problem is ill-posed at the level of device identity.

\subsection{Chains and mixes}

Hinrichs et al.\ estimated effect classes and parameters with the guitar
still inside a mix with bass, keyboards and drums: 97.4\% accuracy on isolated
notes and 90.0\% in mixes, with parameter errors below 0.016 for distortion
and 0.05--0.1 for chains in a fixed order~\cite{hinrichs2022convolutional}.
Guo and McFee recognised combinations of 13 SoX effects on guitar phrases,
with micro and macro F1 of 0.876 and 0.906 on an unseen dataset without the
clean signal~\cite{guo2023automatic}; overdrive and distortion, slapback and
feedback delay, and the modulation effects were the usual confusions. Order is
harder than presence: Wada et al.\ reach a set F1 of 0.987 on three-effect
chains but a macro F1 of 0.756 on their 16 orders~\cite{wada2025hyperbolic},
and Rossi et al.\ an exact-chain accuracy of 0.815 when order and parameters
vary~\cite{rossi2025automatic}. Take et al.\ estimated the chain while
recovering the dry signal by peeling effects one at a time; re-applying the
estimated chain to the recovered dry signal improved the wet reconstruction by
only 0.67\,dB SI-SDR, against 2.24\,dB from the true dry signal, because
removal and estimation errors add up~\cite{take2024audio}.

Okita and Katayose combined prediction and search: a network predicts the dry
signal and the effect types, then CMA-ES or TPE searches the order and the
parameters of a black-box chain to maximise the SI-SDR of the re-rendered wet
signal~\cite{okita2026audio}. With three effects (chorus, distortion,
reverb) the search raised the wet reconstruction from 18.18 to 23.07\,dB and
lowered MR-STFT from 0.465 to 0.340. The chorus rate had to be fixed because
its search landscape was too rugged, a warning for every LFO-driven effect.
Table~\ref{tab:recognition} summarises these results.

\begin{table}[t]
\centering
\caption{Effect recognition and parameter estimation on guitar, as published.
All training data are synthetic renders.}
\label{tab:recognition}
\scriptsize
\setlength{\tabcolsep}{2.5pt}
\begin{tabular}{@{}llll@{}}
\toprule
Work & Input & Task & Result \\
\midrule
Stein~\cite{stein2010dafx} & notes & 6 cascaded effects & 99.2\% acc. \\
Comunità~\cite{comunita2021guitar} & notes & 13 drives, gain, tone & 81--91\%, MAE 0.04 \\
Hinrichs~\cite{hinrichs2022convolutional} & notes in mix & class, params & 90.0\%, MAE $<$0.016 \\
Guo~\cite{guo2023automatic} & phrases & 13 effects & F1 0.876 / 0.906 \\
Rossi~\cite{rossi2025automatic} & DI phrases & 5-effect chain & exact 0.815 \\
Wada~\cite{wada2025hyperbolic} & phrases & ordered chain & set 0.987, order 0.756 \\
Okita~\cite{okita2026audio} & phrases & type, search & 23.07\,dB SI-SDR \\
StemFX~\cite{cheng2026stemfx} & sep.\ stems & FX chain & top-1 86.8\% \\
\bottomrule
\end{tabular}
\end{table}

\subsection{Blind estimation of processing graphs}

Outside guitar, Lee et al.\ estimate a whole processing graph, with nodes,
connections and parameters, from a processed reference with a CNN encoder and
a graph transformer decoder~\cite{lee2023blind}; most errors were wrong orders
or processors hidden by a later destructive one. Their later work fits an
over-complete differentiable console to dry tracks and the mix and prunes
two-thirds of the processors without losing quality~\cite{lee2024searching,lee2025reverse},
using the GRAFX library~\cite{lee2024grafx}. Colonel and Reiss recovered gain,
pan, equalisation, delay and reverb of a mix from its raw tracks with
differentiable processing~\cite{colonel2021reverse}. These methods need the
dry tracks, but the idea of fitting a fixed, over-complete chain and letting
the fit switch parts off transfers to a fixed pedal chain. StemFX is the first
system trained to predict effect chains directly on separated stems without a
dry reference: a band-split encoder and a transformer decoder predict
tokenised chains drawn from 85 effects, trained on pseudo-stems of about 105k
songs~\cite{cheng2026stemfx}. It predicts mixing effects relative to an
original stem, not guitar rigs, and guitar sits in its ``other'' stem.
LLM2Fx-Tools predicts executable effect chains with tool calls from dry and
processed pairs~\cite{doh2026llm2fxtools}.

\subsection{Sub-problems of a guitar chain}

\emph{Modulation.} Mitcheltree et al.\ extract the LFO of phasers, flangers
and choruses from processed audio with 1.8--4.9\% error on guitar and 3.1\% on
unseen real guitar, then use it to model unknown effects~\cite{mitcheltree2023modulation}.
Carson et al.\ fit differentiable flangers, choruses and phasers, and need a
low-frequency weighting of the loss to avoid local minima in delay
time~\cite{carson2026gradient}.

\emph{Delay and reverberation.} Delay time and mix were estimated only on
synthetic notes~\cite{jurgens2020recognizing,hinrichs2022convolutional}.
Blind estimation of artificial reverberation parameters exists for speech with
differentiable reverberators~\cite{lee2022differentiable}, and blind RT60
estimation is a speech and room-acoustics topic~\cite{gamper2018blind}; we
found nothing on the plate, spring or hall reverb of a guitar track inside a
mix.

\emph{Dynamics.} Compressor parameters can be estimated by hybrid neural
regression with model-based inversion~\cite{sun2024neural} or gradient-free
search in a perceptual feature space~\cite{sun2026blackbox}.

\emph{The instrument.} Pickup position and plucking point can be estimated from
the spectrum, also through amplifier chains (median errors of 2--7\,mm for the
pickup)~\cite{mohamad2017pickup,mohamad2017estimating}. The guitar and its
pickups are part of the tone that an effects application cannot change, only
compensate with equalisation.
